\section{Introduction}

Public event logs can reveal every recorded participation while suppressing
the identifier that joins participations into a shared event. A shared-trip
record, for example, may expose departure time, duration, and destination but
omit its vehicle-run key. The analyst can count trips directly; measuring
within-event composition or comparing outcomes across possible co-participants
requires a relation that the release does not supply. A reproducible point
reconstruction selects one explanation of these rows, but its aggregate answer
need not survive other explanations allowed by the same observations.

Aggregate bounds under uncertain linkage are an established response to this
problem \citep{turkcapar2023aggregate,hua2012aggregate}. Our focus is the
additional structure of \emph{temporal events}: participants enter and leave,
capacity limits simultaneous occupancy rather than total membership, and
released timestamps can themselves be uncertain. A two-seat event can contain
three sequential participants even when the first and last never coexist.
Conversely, two individually possible overlaps may require incompatible
realizations of a shared participant's time. Pairwise candidate compatibility
therefore does not establish joint temporal-event feasibility.

The resulting question is specific:
\begin{quote}
Which aggregate conclusions hold across all event partitions that admit one
common timestamp completion satisfying the declared connectivity and
simultaneous-capacity constraints?
\end{quote}
We call the input \emph{relation-incomplete}; row completeness, when assumed,
is relative to a declared candidate universe and observation protocol. It is
not a guarantee that every operational participant appears in a public extract.

\methodname{} represents a feasible world by a timestamp completion and a
partition of selected rows into positive-overlap-connected events. Required
core rows define the cohort, while optional boundary rows can complete its
events without being forced into the cohort themselves. Aggregate extrema
describe the conclusions supported by this model. They are conditional
identified bounds, not confidence intervals or recovered event identities.
The model deliberately omits unobserved routing and provider rules; a
temporally feasible world need not be operationally realizable in every domain.

The computational challenge is the implicit family of possible events.
Even with exact timestamps, optimizing each event independently is insufficient:
events compete for the same rows. We exploit a local interval structure and
retain an integer search for their global coupling. Set-valued timestamps are
handled separately through joint feasibility, rather than assuming that a
polynomial exact-time subproblem also solves uncertain-time instances.

Our contributions are threefold.
\begin{enumerate}[leftmargin=*,itemsep=2pt,topsep=3pt]
\item \textbf{A precise event-level identification model.} We combine required
and optional rows, sequential membership, simultaneous occupancy, and joint
timestamp completion. Separation examples show why event cardinality, pairwise
compatibility, and outer time envelopes cannot substitute for these conditions.
Support nesting and threshold certification specify the resulting guarantee;
they are consequences of possible-world semantics, not new general principles.
\item \textbf{An integral local optimization primitive.} With exact intervals,
a fixed root and span, and additive row weights, interleaved segment-coverage
and boundary-bridge constraints preserve the consecutive-ones structure. The
resulting integral LP prices a single event. We integrate this primitive into
Dantzig--Wolfe decomposition and branch-and-price for support maximization,
with a nonintegral master and explicit bounds when search is incomplete.
\item \textbf{An audit separating three evidentiary questions.} Controlled
ordered-event instances check optimization and decision logic. Annotated
pedestrian dyads test whether geometric candidate rules retain a real relation.
Public Chicago and NYC records test conditional ambiguity and computational
coverage. No experiment is used as evidence for a different layer's claim.
\end{enumerate}

In 3,000 controlled instances, a feasible temporal point rule has
17.4--19.0\% threshold error; the frontier flags all observed errors when truth
is admissible. This is a validation implication of coverage. More informative
is its failure boundary: candidate truncation retains about 84\% of true
members but only 31--33\% of complete true worlds. In the annotated dyad
audit, a pilot-selected support retains 81/82 follow-up worlds yet certifies
only 29/246 decisions. Coverage, sparsity, and decision informativeness are
different properties.

The public results also delimit the contribution. The completed Chicago
24-window pilot has 23 eligible windows and 2,744 certified numerical endpoint
pairs, with 686 pairs excluded for missing public values. Its 96-window
follow-up is incomplete. NYC has 125 witness-certified ambiguous decisions
among 126 outcome--capacity cells, while a separate 18-cell support-maximization
lattice closes through 16 cores and 48 buffers. A fixed-input structural audit
finds \emph{no} endpoint change in 24 comparable restricted-versus-ordered
comparisons. Thus public-data evidence currently supports conditional ambiguity
and bounded-scale computation; it does not yet establish an empirical advantage
from the general event model. This distinction determines both what the paper
claims and what further evidence is needed.
